\section{Výběrový kvantil a jeho asymptotické vlastnosti}

Nechť $F(y)$ je distribuční funkce náhodné veličiny $Y$. Populační kvantil řádu $p$, kde $p \in (0, 1)$, značíme $Q_p$ a definujeme jako:
\begin{equation}
Q_p = F^{-1}(p) = \inf \{ y : F(y) \ge p \}
\end{equation}
Výběrový kvantil $\hat{Q}_p$ z prostého náhodného výběru o rozsahu $n$ je definován jako inverzní funkce k empirické distribuční funkci $\hat{F}_n(y)$.

Klíčovým poznatkem pro statistickou inferenci je asymptotická normalita výběrového kvantilu. Podle Bahadurovy reprezentace pro $n \to \infty$ platí:
\begin{equation}
\sqrt{n}(\hat{Q}_p - Q_p) \xrightarrow{d} N\left( 0, \frac{p(1-p)}{[f(Q_p)]^2} \right)
\end{equation}
Z tohoto vztahu vyplývá asymptotický rozptyl:
\begin{equation}
D(\hat{Q}_p) \approx \frac{p(1-p)}{n [f(Q_p)]^2}
\end{equation}
Zde je jasně patrný problém závislosti na $1/f(Q_p)$. Pokud je hustota v bodě kvantilu malá (chvosty rozdělení), rozptyl je velký.

\section{Taylorova linearizace}

Linearizace je obecná technika pro aproximaci rozptylu nelineární statistiky $g(X)$. Rozvojem Taylorovy řady prvního řádu kolem střední hodnoty $\mu$ dostáváme $g(X) \approx g(\mu) + g'(\mu)(X-\mu)$. Aplikací operátoru rozptylu: $D(g(X)) \approx [g'(\mu)]^2 D(X)$.

V případě kvantilu narážíme na problém, že $\hat{Q}_p$ není jednoduchá funkce průměrů. Můžeme však linearizovat vztah mezi kvantilem a distribuční funkcí. Vyjdeme z rovnosti $\hat{F}(\hat{Q}_p) \approx p$. Taylorův rozvoj funkce $\hat{F}$ v bodě $\hat{Q}_p$ kolem $Q_p$ dává:
\begin{equation}
\hat{F}(\hat{Q}_p) \approx \hat{F}(Q_p) + (\hat{Q}_p - Q_p) f(Q_p)
\end{equation}
Odtud vyjádříme chybu odhadu:
\begin{equation}
(\hat{Q}_p - Q_p) \approx \frac{p - \hat{F}(Q_p)}{f(Q_p)}
\end{equation}
Tento vztah převádí "vodorovnou" odchylku kvantilů na "svislou" odchylku distribuční funkce, škálovanou sklonem tečny (hustotou).

\section{Woodruffova metoda}

Woodruff (1952) navrhl elegantní metodu, která využívá výše uvedený princip linearizace, ale obchází nutnost znát $f(Q_p)$.

Postup sestává z následuících kroků:
\begin{enumerate}
    \item Sestrojíme interval spolehlivosti pro očekávanou hodnotu distribuční funkce, tedy pro $p$. Pro hladinu významnosti $\alpha$:
    \begin{equation}
    p_L = p - u_{1-\alpha/2} \sqrt{\frac{p(1-p)}{n}}, \quad p_U = p + u_{1-\alpha/2} \sqrt{\frac{p(1-p)}{n}}
    \end{equation}
    \item Najdeme výběrové kvantily odpovídající těmto mezním pravděpodobnostem. Získáme tak interval $[ \hat{Q}_{p_L}, \hat{Q}_{p_U} ]$. Tento interval je $100(1-\alpha)\%$ intervalem spolehlivosti pro $Q_p$.
    \item Pokud potřebujeme odhad směrodatné odchylky (Standard Error), odvodíme ji z délky intervalu (za předpokladu normality):
    \begin{equation}
    \hat{SE}(\hat{Q}_p) = \frac{\hat{Q}_{p_U} - \hat{Q}_{p_L}}{2 u_{1-\alpha/2}}
    \end{equation}
\end{enumerate}
Tato metoda implicitně "odhaduje" převrácenou hodnotu hustoty pomocí šířky intervalu mezi kvantily, což je mnohem stabilnější než odhad hustoty v jednom bodě.
